% ^^A -*- japanese-latex -*-
%
\ifx\epTeXinputencoding\undefined\else
  \epTeXinputencoding utf8 % ^^A added (2017-10-04)
\fi
%
% \iffalse
%
%   gcmcline.dtx
%
%<*driver>
\documentclass[uplatex,a4paper,12Q,dvipdfmx]{jsarticle}
\usepackage{doc}
\usepackage{metalogo}
\providecommand*{\eTeX}{$\varepsilon$-\TeX}
\providecommand*{\pLaTeX}{p\kern-.05em\LaTeX}
\providecommand*{\pLaTeXe}{p\kern-.05em\LaTeXe}
\providecommand*{\upTeX}{u\pTeX}
\providecommand*{\XeTeX}{XeTeX}
\usepackage{ifptex}
\ifptex
  \xspcode"5C=1 %% \
  \xspcode"22=1 %% "
\fi
\addtolength{\textwidth}{-1.2in}
%% \addtolength{\evensidemargin}{1.5in}
\addtolength{\oddsidemargin}{1.5in}
%% \addtolength{\marginparwidth}{1.5in}
\setlength\marginparpush{0pt}
\usepackage{etoolbox}
\usepackage{listings}
\usepackage{color}
\usepackage{xspace}
\usepackage{ifpdf}
\usepackage{hyperref}
\ifptex
  \usepackage{pxjahyper}
\fi
\usepackage[T1]{fontenc}
\usepackage{lmodern}
\definecolor{myblue}{rgb}{0,0,0.75}
\definecolor{mygreen}{rgb}{0,0.45,0}
\hypersetup{colorlinks,
            hyperfootnotes=false,
            linkcolor=myblue,
            urlcolor=mygreen}
\CodelineNumbered
\EnableCrossrefs
% \CodelineIndex
\RecordChanges                  % Gather update information
\setcounter{StandardModuleDepth}{1}

\usepackage{musixtex}
\usepackage{gcmcline}

\makeatletter

\newrobustcmd*{\file}[1]{\mbox{\ttfamily#1}}
\newrobustcmd*{\pkg}[1]{\mbox{\sffamily#1}}
\newcommand*{\Lcount}[1]{\textsl{\small#1}}
\newcommand*{\Lopt}[1]{\textsf{#1}}
\newcounter{@clineno}
\def\mlineplus#1{\setcounter{@clineno}{\arabic{CodelineNo}}%
   \addtocounter{@clineno}{#1}\arabic{@clineno}}
\def\cmd#1{\cs{\expandafter\cmd@to@cs\string#1}}
\def\cmd@to@cs#1#2{\char\number`#2\relax}
\DeclareRobustCommand{\cs}[1]{\texttt{\char`\\#1}}
\DeclareRobustCommand{\env}[1]{\texttt{#1}}
\providecommand{\marg}[1]{%
  {\ttfamily\char`\{}\meta{#1}{\ttfamily\char`\}}}
\providecommand{\oarg}[1]{%
  {\ttfamily[}\meta{#1}{\ttfamily]}}
\def\TODO{\par\textsf{To Do:\ }}%
\makeatother

\GetFileInfo{gcmcline.dtx}
\begin{document}
  \DocInput{gcmcline.dtx}
\end{document}
%</driver>
%
% \fi
% \makeatletter
% \IndexPrologue{%
%   \section*{索引}%
%   \markboth{索引}{索引}%
%   イタリック体の数字はその項目の記述されたページへの，
%   下線の引かれた数字はその項目の定義された
%   \ifcodeline@index
%     コードの行
%   \else
%     ページ
%   \fi
%   への，ローマン体の数字はその項目の使われた
%   \ifcodeline@index
%     コードの行
%   \else
%     ページ
%   \fi
%   への参照となっている．
% }
% \GlossaryPrologue{%
%   \section*{{変更履歴}}%
%   \markboth{{変更履歴}}{{変更履歴}}%
% }
% \makeatother
% \changes{v1.0.0}{2025/11/28}
%   {make this .dtx file}
% \iffalse TeX commands \fi
% \DoNotIndex{\ ,\-,\/}
% \DoNotIndex{\above,\abovedisplayshortskip,\abovedisplayskip}
% \DoNotIndex{\abovewithdelims,\accent,\adjdemerits}
% \DoNotIndex{\advance,\afterassignment,\aftergroup}
% \DoNotIndex{\atop,\atopwithdelims}
% \DoNotIndex{\badness,\baselineskip,\batchmode}
% \DoNotIndex{\begingroup,\belowdisplayshortskip,\belowdisplayskip}
% \DoNotIndex{\binoppenalty,\botmark,\box}
% \DoNotIndex{\boxmaxdepth,\brokenpenalty}
% \DoNotIndex{\catcode,\char,\chardef}
% \DoNotIndex{\cleaders,\closein,\closeout}
% \DoNotIndex{\clubpenalty,\copy,\count}
% \DoNotIndex{\countdef,\cr,\crcr,\csname}
% \DoNotIndex{\day,\cr,\crcr,\csname}
% \DoNotIndex{\countdef,\deadcycles,\def,\defaulthyphenchar}
% \DoNotIndex{\defaultskewchar,\delcode,\delimiter,\delimiterfactor}
% \DoNotIndex{\delimitershortfall,\dimen,\dimendef,\discretionary}
% \DoNotIndex{\displayindent,\displaylimits,\displaystyle,\displaywidowpenalty}
% \DoNotIndex{\displaywidth,\divide,\doublehyphendemerits,\dp,\dump}
% \DoNotIndex{\edef,\else,\emergencystretch,\end}
% \DoNotIndex{\endcsname,\endgroup,\endinput,\endlinechar}
% \DoNotIndex{\eqno,\errhelp,\errmessage,\errorcontextlines}
% \DoNotIndex{\errorstopmode,\escapechar,\everycr,\everydisplay}
% \DoNotIndex{\everyhbox,\everyjob,\everymath,\everypar}
% \DoNotIndex{\everyvbox,\exhyphenpenalty,\expandafter}
% \DoNotIndex{\fam,\fi,\finalhyphendemerits,\firstmark}
% \DoNotIndex{\floatingpenalty,\font,\fontdimen,\fontname,\futurelet}
% \DoNotIndex{\gdef,\global,\globaldefs}
% \DoNotIndex{\halign,\hangafter,\hangindent,\hbadness}
% \DoNotIndex{\hbox,\hfil,\hfill,\hfilneg}
% \DoNotIndex{\hfuzz,\hoffset,\holdinginserts,\hrule}
% \DoNotIndex{\hsize,\hskip,\hss,\ht}
% \DoNotIndex{\hyphenation,\hyphenchar,\hyphenpenalty}
% \DoNotIndex{\if,\ifcase,\ifcat,\ifdim}
% \DoNotIndex{\ifeof,\iffalse,\ifhbox,\ifhmode}
% \DoNotIndex{\ifinner,\ifmmode,\ifnum,\ifodd}
% \DoNotIndex{\iftrue,\ifvbox,\ifvmode,\ifvoid}
% \DoNotIndex{\ifx,\ignorespaces,\immediate,\indent}
% \DoNotIndex{\input,\inputlineno,\insert,\insertpenalties,\interlinepenalty}
% \DoNotIndex{\jobname}
% \DoNotIndex{\kern}
% \DoNotIndex{\language,\lastbox,\lastkern,\lastpenalty}
% \DoNotIndex{\lastskip,\lccode,\leaders,\left}
% \DoNotIndex{\lefthyphenmin,\leftskip,\leqno,\let}
% \DoNotIndex{\limits,\linepenalty,\lineskip,\lineskiplimit}
% \DoNotIndex{\long,\looseness,\lower,\lowercase}
% \DoNotIndex{\mag,\mark,\mathaccent,\mathbin}
% \DoNotIndex{\mathchar,\mathchardef,\mathchoice,\mathclose}
% \DoNotIndex{\mathcode,\mathinner,\mathop,\mathopen}
% \DoNotIndex{\mathord,\mathpunct,\mathrel,\mathsurround}
% \DoNotIndex{\maxdeadcycles,\maxdepth,\meaning,\medmuskip}
% \DoNotIndex{\message,\mkern,\month,\moveleft}
% \DoNotIndex{\moveright,\mskip,\multiply,\muskip,\muskipdef}
% \DoNotIndex{\newlinechar,\noalign,\noboundary,\noexpand}
% \DoNotIndex{\noindent,\nolimits,\nonscript,\nonstopmode}
% \DoNotIndex{\nulldelimiterspace,\nullfont,\number}
% \DoNotIndex{\omit,\openin,\openout,\or}
% \DoNotIndex{\outer,\output,\outputpenalty,\over}
% \DoNotIndex{\overfullrule,\overline,\overwithdelims}
% \DoNotIndex{\pagedepth,\pagefilllstretch,\pagefillstretch,\pagefilstretch}
% \DoNotIndex{\pagegoal,\pageshrink,\pagestretch,\pagetotal}
% \DoNotIndex{\par,\parfillskip,\parindent,\parshape}
% \DoNotIndex{\parskip,\patterns,\pausing,\penalty}
% \DoNotIndex{\postdisplaypenalty,\predisplaypenalty,\predisplaysize}
% \DoNotIndex{\pretolerance,\prevdepth,\prevgraf}
% \DoNotIndex{\radical,\raise,\read,\relax}
% \DoNotIndex{\relpenalty,\right,\righthyphenmin,\rightskip,\romannumeral}
% \DoNotIndex{\scriptfont,\scriptscriptfont,\scriptscriptstyle,\scriptspace}
% \DoNotIndex{\scriptstyle,\scrollmode,\setbox,\setlanguage}
% \DoNotIndex{\sfcode,\shipout,\show,\showbox}
% \DoNotIndex{\showboxbreadth,\showboxdepth,\showlists,\showthe}
% \DoNotIndex{\skewchar,\skip,\skipdef,\spacefactor}
% \DoNotIndex{\spaceskip,\span,\special,\splitbotmark}
% \DoNotIndex{\splitfirstmark,\splitmaxdepth,\splittopskip,\string}
% \DoNotIndex{\tabskip,\textfont,\textstyle,\the}
% \DoNotIndex{\thickmuskip,\thinmuskip,\time,\toks}
% \DoNotIndex{\toksdef,\tolerance,\topmark,\topskip}
% \DoNotIndex{\tracingcommands,\tracinglostchars,\tracingmacros,\tracingonline}
% \DoNotIndex{\tracingoutput,\tracingpages,\tracingparagraphs,\tracingrestores}
% \DoNotIndex{\tracingstats}
% \DoNotIndex{\uccode,\uchyph,\underline,\unhbox}
% \DoNotIndex{\unhcopy,\unkern,\unpenalty,\unskip}
% \DoNotIndex{\unvbox,\unvcopy,\uppercase}
% \DoNotIndex{\vadjust,\valign,\vbadness,\vbox}
% \DoNotIndex{\vcenter,\vfil,\vfill,\vfilneg}
% \DoNotIndex{\vfuzz,\voffset,\vrule,\vsize}
% \DoNotIndex{\vskip,\vsplit,\vss,\vtop}
% \DoNotIndex{\wd,\widowpenalty,\write}
% \DoNotIndex{\xdef,\xleaders,\xspaceskip}
% \DoNotIndex{\year}

% \iffalse e-TeX commands \fi
% \DoNotIndex{\protected,\detokenize,\unexpanded}
% \DoNotIndex{\readline,\scantokens}
% \DoNotIndex{\eTeXrevision,\eTeXversion,\currentgrouplevel,\currentgrouptype}
% \DoNotIndex{\ifcsname,\ifdefined,\lastnodetype}
% \DoNotIndex{\marks}
% \DoNotIndex{\beginL,\beginR,\endL,\endR}
% \DoNotIndex{\TeXXeTstate,\predisplaydirection}
% \DoNotIndex{\interactionmode,\showgroups,\showtokens}
% \DoNotIndex{\tracingscantokens,\tracinggroups}
% \DoNotIndex{\everyeof,\middle,\unless}
%
% \title{gcmcline}
% \author{bell}
% \date{\filedate}
% \maketitle
% \tableofcontents
%
% \MakeShortVerb{\|}
%
% \section{概説}
% 楽譜上での直線を実装します．
% 
% \DescribeMacro{\imusline}\cmd{\imusline}は
% \cmd{\imusline}\oarg{text}\marg{id}\marg{height}の形で直線を開始します．
% \begin{itemize}
% \item \meta{text}はmuslineの上に出力される文字を設定します．
% \item \meta{id}は参照するmusline IDを設定します．
% \item \meta{height}はmuslineの出力される高さを設定します．
% \end{itemize}
%
% このようにして開始した直線は
% その地点の右側1.5\cmd{\qn@width}のところを始点とします．
%
% \DescribeMacro{\tmusline}\cmd{\tmusline}は
% \cmd{\tmusline}\marg{id}\marg{height}の形で直線を終了します．
% \begin{itemize}
% \item \meta{id}は参照するmusline IDを設定します．
% \item \meta{height}はmuslineの出力される終端部の高さを設定します．
% \end{itemize}
%
% \StopEventually{}
% \section{実装}
%
%    \begin{macrocode}
%<*macro>
%%=============================================================================%
%%                                                                             %
%%   GGGGG          A          CCCCC    H       H    IIIII                     %
%%  G     G        A A        C     C   H       H      I                       %
%% G       G      A   A      C       C  H       H      I                       %
%% G             A     A     C          H       H      I                       %
%% G            A       A    C          HHHHHHHHH      I     --------          %
%% G     GGGG   AAAAAAAAA    C          H       H      I                       %
%% G       G    A       A    C       C  H       H      I                       %
%%  G     GG    A       A     C     C   H       H      I                       %
%%   GGGGG G    A       A      CCCCC    H       H    IIIII                     %
%%                                                                             %
%%                M         M   UUU     UUU    CCCCC    H       H    IIIII     %
%%                MM       MM    U       U    C     C   H       H      I       %
%%                M M     M M    U       U   C       C  H       H      I       %
%%                M M     M M    U       U   C          H       H      I       %
%%                M  M   M  M    U       U   C          HHHHHHHHH      I       %
%%                M  M   M  M    U       U   C          H       H      I       %
%%                M   M M   M    U       U   C       C  H       H      I       %
%%                M   M M   M     U     U     C     C   H       H      I       %
%%                M    M    M      UUUUU       CCCCC    H       H    IIIII     %
%%                                                                             %
%%=============================================================================%
\NeedsTeXFormat{LaTeX2e}% 
%    \end{macrocode}
% \begin{macro}{\gcmc@line@Version@Major}
% このパッケージのメジャーバージョン番号です．
%    \begin{macrocode}
\newcommand{\gcmc@line@Version@Major}{1}
%    \end{macrocode}
% \end{macro}
% \begin{macro}{\gcmc@line@Version@Minor}
% このパッケージのマイナーバージョン番号です．
%    \begin{macrocode}
\newcommand{\gcmc@line@Version@Minor}{0}
%    \end{macrocode}
% \end{macro}
% \begin{macro}{\gcmc@line@Version@Patch}
% このパッケージのパッチバージョン番号です．
%    \begin{macrocode}
\newcommand{\gcmc@line@Version@Patch}{0}% 2025-11-28
%    \end{macrocode}
% \end{macro}
% \begin{macro}{\gachimuchimacro@Version}
% このパッケージのバージョン番号です．
%    \begin{macrocode}
\newcommand{\gcmc@line@Version}{%
  \gcmc@line@Version@Major
  .\gcmc@line@Version@Minor
  .\gcmc@line@Version@Patch
}
%    \end{macrocode}
% \end{macro}
%    \begin{macrocode}
\ProvidesPackage{gcmcline}
  [2025/11/28 Gachimuchi Koza Macro set ver.~\gcmc@line@Version]
%    \end{macrocode}
% \subsection{muslineの追加}
% \paragraph{概説}
% スラーと同様の方法で直線を使えるようにします．
% \paragraph{実装}
% \subparagraph{下準備}
% \begin{macro}{\N@musline}
% 保留された直線の数を保持するカウンタ・レジスタを用意します．
%    \begin{macrocode}
%% muslineの追加
%% registers for musline
\newcount\N@musline
%    \end{macrocode}
% \end{macro}
% \begin{macro}{\test@muslinenum}
% 参照すべきmusline IDが範囲外であるときに
% エラー・メッセージを出す用のテストマクロです．
%    \begin{macrocode}
\def\test@muslinenum{%
  \ifnum\n@i<\z@ \n@i\@c \fi
  \ifnum\n@i<\maxmuslines \else
    \count@\maxmuslines \advance\count@\m@ne
    \errmessage{Wrong musline reference number \the\n@i!
                (valid: 0 to \the\count@)}%
    \n@i\z@
  \fi
  \advance\n@i\@ne}%

%    \end{macrocode}
% \end{macro}
% \begin{macro}{\@oldmaxmuslines}
% 現在の最大musline IDを保持します．
%    \begin{macrocode}
\chardef\@oldmaxmuslines=\z@
%    \end{macrocode}
% \end{macro}
% \begin{macro}{\setmaxmuslines}\marg{max-ID number}\par
% \meta{max-ID number}でmusline IDの最大値を設定し，
% 各種パラメータを定義します．
% \begin{itemize}
% \item \cmd{\musline@x}: 始点でのmuslineの水平位置[dimen]
% \item \cmd{\musline@s}: muslineの向き(u, d)およびmuslineが使用中ではないことを意味する
%                 フラグ(x)
% \item \cmd{\musline@Y}: muslineの垂直方向のオフセット
% \item \cmd{\musline@a}: musline開始時の\cmd{\altportee}
% \item \cmd{\musline@y}: 始点でのmuslineの垂直位置[steps of internote]
% \item \cmd{\musline@N}: muslineの使われている楽器の番号
% \item \cmd{\musline@z}: (予定)muslineが\cmd{\breakmusline}により中断されたときに
%                         指定する終了時の高さ[steps of internote]
% \item \cmd{\musline@t}: muslineの上または下に表示されるテキスト・トークン
\iffalse
% \item \cmd{\musline@V}: 終点でのスラーの高さのオフセット；
%                 それに伴い\cmd{\s@Y}は始点での垂直方向の
%                 オフセットとして扱われる.
\fi
% \item \cmd{\musline@o}: muslineの鉤の開き．
% \end{itemize}
%    \begin{macrocode}
\def\setmaxmuslines#1{%
  \ifnum#1>\@oldmaxmuslines
    \chardef\maxmuslines=#1\relax
    \ifnum#1>\@c
      \errmessage{The maximum reference number of elements
                  is limited to \the\@c}
      \chardef\maxmuslines=\@c
    \fi
    \count@\@oldmaxmuslines\loop
      \advance\count@ by\@ne
      % horizontal start position of musline
      \expandafter \noexpand@newdimsk \csname musline@x\roman@c@\endcsname
      %
      % sense u,d and x (flag) for pending musline
      \expandafter \noexpand@newtoks  \csname musline@s\roman@c@\endcsname
      \csname musline@s\roman@c@\endcsname={x}%
      %
      % vertical offset of musline
      \expandafter \noexpand@newdimsk \csname musline@Y\roman@c@\endcsname
      %
      % altportee of musline
      \expandafter \noexpand@newskip  \csname musline@a\roman@c@\endcsname
      %
      % vertical start position of slur [steps of internote]
      \expandafter \noexpand@newcount \csname musline@y\roman@c@\endcsname
      %
      % instrument number of that slur (to retrieve internote)
      \expandafter \noexpand@newcount \csname musline@N\roman@c@\endcsname
      %
      % vertical end position of slur [steps of internote], used from
      % \breakslur.  the value \maxdimen is used to decide, if
      % \breakslur is used or not (flag)
      \expandafter \noexpand@newcount \csname musline@z\roman@c@\endcsname
      \csname musline@z\roman@c@\endcsname\maxdimen
      %
      % text token
      \expandafter \noexpand@newtoks  \csname musline@t\roman@c@\endcsname
      %
      %
      % open left/right token
      \expandafter \noexpand@newtoks  \csname musline@o\roman@c@\endcsname
      \csname musline@o\roman@c@\endcsname={x}%
      %
    \ifnum\count@<\maxmuslines\repeat
    \expandafter\message\expandafter{maxmuslines=\the\maxmuslines}%
    \chardef\@oldmaxmuslines=\maxmuslines\relax%
    \normalnotesize\resetlayout
  \fi}%

%    \end{macrocode}
% \end{macro}
% \begin{macro}{\s@l@ctmusline}\marg{musline ID}\par
% IDが\meta{musline ID}であるmuslineのパラメータを取得します．
%    \begin{macrocode}
\def\s@l@ctmusline#1\relax{%
  \n@i#1\relax\test@muslinenum%
  \edef\musline@a{\csname musline@a\roman@n@i}%  dimen (altportee)
  \edef\musline@x{\csname musline@x\roman@n@i}%  dimen (hor start pos)
  \edef\musline@y{\csname musline@y\roman@n@i}%  count (startnote [internote])
  \edef\musline@N{\csname musline@N\roman@n@i}%  count (instrument number)
  \edef\musline@z{\csname musline@z\roman@n@i}%  count (endnote [internote])
  \edef\musline@Y{\csname musline@Y\roman@n@i}%  dimen (voffset)
  \edef\musline@t{\csname musline@t\roman@n@i}%  token (text)
  \edef\musline@s{\csname musline@s\roman@n@i}%  token (sense)
  \edef\musline@o{\csname musline@o\roman@n@i}%  token (sense)
}

%    \end{macrocode}
% \end{macro}
% \subparagraph{始点}
% \begin{macro}{\imusline}\oarg{text}\marg{id}\marg{height}\par
% muslineの始点を設定します．
% \begin{itemize}
% \item \meta{text}はmuslineの上に出力される文字を設定します．
% \item \meta{id}は参照するmusline IDを設定します．
% \item \meta{height}はmuslineの出力される高さを設定します．
% \end{itemize}
%    \begin{macrocode}
%% \ibktu/d #1=text(optional) #2=musline id #3=height[\internote]
\newcommand{\imusline}[1][]{\@i@musline[#1]{1.5}{\z@}u}
%    \end{macrocode}
% \end{macro}
% \begin{macro}{\@i@musline}\oarg{text}\marg{hoffset}\marg{voffset}%
%              \marg{sence}\marg{id}\marg{height}\par
% muslineの始点を設定する内部マクロです．
% \begin{itemize}
% \item \meta{text}はmuslineの上または下に出力される文字を設定します．
% \item \meta{hoffset}は現在のカーソルから右へずらす長さの係数を設定します．
%       単位は符頭の長さ\cmd{\qn@width}です．
% \item \meta{voffset}は始点の高さから上へずらす長さを設定します．
% \item \meta{sence}はmuslineの上下の向きを設定します．
% \item \meta{id}は参照するmusline IDを設定します．
% \item \meta{height}はmuslineの出力される始点の高さを設定します．
% \end{itemize}
% 設定されたパラメータは\cmd{\musline@}から始まる
% 各マクロあるいはレジスタに（グローバルに）定義あるいは代入されます．
%    \begin{macrocode}
%% \@i@musline #1=text #2=hoffset[\qn@width] #3=voffset
%%         #4=musline sence(u,d) #5=musline id #6=height[\internote]
\def\@i@musline[#1]#2#3#4#5#6{%
  \check@staff
    \global\advance\N@musline\@ne % update musline counter
    \s@l@ctmusline#5\relax
    % test for already invoked \islur
    \if x\the\musline@s \else\errmessage{\@mis\noexpand\tmusline#5}\fi
    \global\musline@Y#3%            store voffset (abs. dim. to rel. height)
    \global\musline@s{#4}%          store sense (u,d)
    \global\musline@t{#1}%          store text
    \global\musline@N\noinstrum@nt% store instrument number
    \inhgetn@i#6\relax
    \global\musline@y\n@i %         start vpos (steps of internote)
    \global\musline@a\altportee %   store altportee of current slur
    \getcurpos
    \advance\y@v#2\qn@width
    \global\musline@x\y@v %         start hpos
  \fi% start hpos (current pos including hoffset)
}

%    \end{macrocode}
% \end{macro}
% \subparagraph{終点}
% \begin{macro}{\tmusline}\marg{id}\marg{height}\par
% muslineの終点を設定します．
% \begin{itemize}
% \item \meta{id}は参照するmusline IDを設定します．
% \item \meta{height}はmuslineの出力される終端部の高さを設定します．
% \end{itemize}
%    \begin{macrocode}
\def\tmusline{\t@musline0}
%    \end{macrocode}
% \end{macro}
% \begin{macro}{\t@musline}\marg{hoffset}\marg{id}\marg{height}\par
%    \begin{macrocode}
%% \t@musline #1=hoffset #2=reference number #3=end note
\def\t@musline#1#2#3{%
  \check@staff
    \s@l@ctmusline#2\relax
    \y@iv#1\qn@width
    \def\@sense{\the\musline@s}%
% test for missing \ibkt
    \if x\@sense \errmessage{\@mis\noexpand\imusline#2}\fi
% compute length
    \getcurpos
    \advance\y@v\y@iv
% eoline
    \y@eol\advance\y@-\beforeruleskip
% clip musline at eoline
    \ifdim\y@v>\y@ \y@v\y@ \advance\y@v\beforeruleskip \y@iv\beforeruleskip \fi
    \advance\y@v-\musline@x
    \if r\the\musline@o \musline@o{l}\fi
    \inhgetn@i#3\relax
    \write@musline{\musline@y}{\n@i}{\y@v}{-\y@iv}%
% reset sense of musline
    \global\musline@s{x}\let\T@i\empty
    \global\musline@o{x}\let\T@i\empty
    \global\advance\N@musline\m@ne% update slur counter
  \fi
}

%    \end{macrocode}
% \end{macro}
% \begin{macro}{\write@musline}
% muslineの描画部分です．
%    \begin{macrocode}
% #1 start pitch #2 end pitch #3 length #4 right hoffset
\def\write@musline#1#2#3#4{\check@staff %
  \ifx\musline@N\undefined
    \relax %
  \else
    \edef\internote{\csname i@n\romannumeral\musline@N\endcsname}%
  \fi
% \immediate\write16{\noexpand\write@musline internote:\the\internote}%
  \n@iv#1%
  \n@vi#2%
  \y@iii#3%
%%
%%   Get note vertical offsets
%%   \n@iv=   vertical offset of first  note(n@i temp);
%%   \n@vi=   vertical offset of second note(n@ii temp);
%%   \y@i =   reference height of first note (pt);
%%   \y@ii=   reference height of end note (pt);
%%   \y@iii=  musline length (pt);
%%
  \n@i\n@vi \pl@base \y@ii\y@i \n@ii\n@vi% \y@ii = end height
  \n@i\n@iv \pl@base % \y@i = initial height
  \advance\y@i\musline@Y \advance\y@ii\musline@Y% add the voffsets
  \ifx\@Ti\@ne
    \advance\y@i\musline@a
    \advance\y@ii\musline@a
    \let\@Ti\empty
  \fi% call from cutslur
%%
%% draw musline
%%
  \llap{%
    \tikz[baseline=0.0ex]{%
      \draw[semithick]
         (0,\y@i)-- (\y@iii,\y@ii);
      \node at ($(\h@lf\y@iii,\y@ii)+(0,\internote)$)
         [above] {\the\musline@t};
    }%
    \hspace{-\y@iv}%
  }%
\fi} % end writ@musline
\setmaxmuslines6

%    \end{macrocode}
% \end{macro}
% \begin{macro}{\C@musline}
% （未実装）シンプルなmusline用マクロです．
% ただし名前被りをどう避けるかが思いついてないので未実装です．
% \end{macro}
%
% 以上で終わりです．
%    \begin{macrocode}
%</macro>
\endinput
%    \end{macrocode}
% \section{To Do}
% \begin{itemize}
% \item
%   senceは存在判定にしか使っていないので整理できるのでは．
% \item
%   開始時・終了時に\meta{dimension}で指定できる垂直・水平オフセットを導入する．
%   これは楽譜の五線ないし加線上から直線が引かれると見た目がよろしくないので，
%   これを避けるためである．
%   \pkg{keyval}式でもいいかも．
% \item
%   \file{gachimuchimacro}内の\cmd{\Mline}等との兼ね合い：引数の形について．
% \end{itemize}
% \Finale
% \PrintIndex \PrintChanges
